\documentclass[a4paper]{article}
% Español (México) con XeLaTeX: fontspec para fuentes Unicode, polyglossia
% para la división silábica y los nombres del documento en la variante mexicana.
\usepackage{fontspec}
\usepackage{polyglossia}
\setdefaultlanguage[variant=mexican]{spanish}
% Los nombres chinos van como «pinyin (original)»: xeCJK da a los caracteres
% Han su propia fuente; sin ella XeLaTeX los omite con un aviso.
% FandolSong no cubre algunos caracteres de nombres (U+73FA); WenQuanYi Zen Hei sí.
\usepackage[AutoFallBack=true]{xeCJK}
\setCJKmainfont{FandolSong-Regular.otf}
\setCJKfallbackfamilyfont{\CJKrmdefault}{WenQuanYi Zen Hei}
% polyglossia no traduce los nombres de listings.
\AtBeginDocument{\providecommand{\lstlistingname}{}\renewcommand{\lstlistingname}{Listado}%
  \providecommand{\lstlistlistingname}{}\renewcommand{\lstlistlistingname}{Índice de listados}}
\usepackage{amsmath,amssymb}
\usepackage{graphicx}
\usepackage[margin=2.35cm]{geometry}
\usepackage[most]{tcolorbox}
\usepackage{etoolbox}
\usepackage{listings}
\usepackage{xcolor}
\usepackage{hyperref}
\usepackage{booktabs,longtable,tabularx,array}
\usepackage{subcaption}
\usepackage{float}
\usepackage{enumitem}
\usepackage{microtype}
\usepackage{fancyhdr}
\usepackage{needspace}

\definecolor{kbBack}{HTML}{F2F6FA}
\definecolor{kbFrame}{HTML}{9DB4CC}
\definecolor{kbTitle}{HTML}{52708F}
\definecolor{ibBack}{HTML}{FBF5E7}
\definecolor{ibFrame}{HTML}{C9A961}
\definecolor{ibTitle}{HTML}{7D5F1E}
\definecolor{wbBack}{HTML}{FCF2F0}
\definecolor{wbFrame}{HTML}{D6A096}
\definecolor{wbTitle}{HTML}{9E4F43}
\definecolor{dbBack}{HTML}{F4F7F3}
\definecolor{dbFrame}{HTML}{AEC2AC}
\definecolor{dbTitle}{HTML}{5F7F64}

\tcbset{softbox/.style={
  enhanced,breakable,boxrule=0.8pt,arc=2pt,left=3mm,right=3mm,
  top=2mm,bottom=2mm,coltitle=white,fonttitle=\bfseries,
  toptitle=0.6mm,bottomtitle=0.6mm
}}
\newtcolorbox{knowledgebox}[1]{softbox,colback=kbBack,colframe=kbFrame,colbacktitle=kbTitle,title={#1}}
\newtcolorbox{importantbox}[1]{softbox,colback=ibBack,colframe=ibFrame,colbacktitle=ibTitle,title={#1}}
\newtcolorbox{warningbox}[1]{softbox,colback=wbBack,colframe=wbFrame,colbacktitle=wbTitle,title={#1}}
\newtcolorbox{dialoguebox}[1]{softbox,colback=dbBack,colframe=dbFrame,colbacktitle=dbTitle,title={#1}}

\lstset{
  language=Python,basicstyle=\ttfamily\small,
  keywordstyle=\color{kbTitle}\bfseries,stringstyle=\color{wbTitle},
  commentstyle=\color{dbTitle},breaklines=true,frame=single,numbers=left,
  numberstyle=\tiny\color{gray},captionpos=b,columns=fullflexible,
  keepspaces=true,showstringspaces=false,extendedchars=true
}
\BeforeBeginEnvironment{lstlisting}{\Needspace{12\baselineskip}}

\hypersetup{colorlinks=true,linkcolor=kbTitle,urlcolor=kbTitle,citecolor=wbTitle}
\setlist{itemsep=0.25em,topsep=0.4em}
\setlength{\parindent}{2em}
\setlength{\parskip}{0.25em}
\newcolumntype{L}[1]{>{\raggedright\arraybackslash}p{#1}}

\providecommand{\notetitle}{Notas del curso: [tema del curso]}
\providecommand{\noteauthors}{Elaboradas a partir de los materiales públicos de Stanford CS25}
\providecommand{\notedate}{Versión reescrita en 2026}
\providecommand{\videochannel}{Stanford Online}
\providecommand{\videopublishdate}{2022}
\providecommand{\videoduration}{[duración del video]}
\providecommand{\videourl}{}
\providecommand{\videocoverpath}{cover.jpg}
\providecommand{\coursesite}{https://web.stanford.edu/class/cs25/}
\providecommand{\repourl}{https://github.com/hqhq1025/ai-course-notes}
\providecommand{\skillversion}{wdkns/wdkns-skills@39f1a04}

\newcommand{\figsource}[1]{\par\vspace{-0.3em}{\footnotesize\color{gray}\emph{Fuente: #1}}\par}
\newcommand{\teachervoice}[1]{\begin{knowledgebox}{Nota de clase}#1\end{knowledgebox}}
\newcommand{\term}[1]{\textbf{#1}}

\pagestyle{fancy}
\fancyhf{}
\fancyfoot[C]{\thepage}
\fancyfoot[R]{\footnotesize\color{gray}\href{\repourl}{AI Course Notes}}
\renewcommand{\headrulewidth}{0pt}
\renewcommand{\footrulewidth}{0pt}

\newcommand{\makecscover}{
\begin{titlepage}
\centering
\vspace*{0.7cm}
{\Large Stanford CS25 · Transformers United\par}
\vspace{0.8cm}
{\huge\bfseries \notetitle\par}
\vspace{0.6cm}
{\large \noteauthors\par}
\vspace{0.25cm}
{\large \notedate\par}
\vspace{0.9cm}
\ifdefempty{\videocoverpath}{}{
  \includegraphics[width=0.86\textwidth,height=0.43\textheight,keepaspectratio]{\videocoverpath}\par
}
\vfill
\begin{tcolorbox}[width=0.92\textwidth,colback=black!2!white,colframe=black!45,boxrule=0.7pt,arc=2pt]
\textbf{Autor o canal del video}: \videochannel\par
\textbf{Fecha de publicación}: \videopublishdate\par
\textbf{Duración del video}: \videoduration\par
\textbf{Sitio del curso}: \href{\coursesite}{\nolinkurl{\coursesite}}\par
\ifdefempty{\videourl}{}{\textbf{Enlace del video}: \href{\videourl}{\nolinkurl{\videourl}}\par}
\textbf{Normas de elaboración}: \skillversion\ + las normas source-first / teacher-voice / visual-QA de este repositorio
\end{tcolorbox}
\end{titlepage}
\tableofcontents
\newpage
}

\usepackage{multicol}

\renewcommand{\notetitle}{Lecture 29: Demystifying Mixtral of Experts}
\renewcommand{\noteauthors}{Elaborado a partir de la clase de Albert Jiang, las slides recuperadas de la grabación oficial y los subtítulos automáticos}
\renewcommand{\notedate}{CS25 V4 · versión reescrita source-first de 2026}
\renewcommand{\videochannel}{Stanford Online}
\renewcommand{\videopublishdate}{Clase: 2024-04-25; publicación: 2024-05-16}
\renewcommand{\videoduration}{1:04:31}
\renewcommand{\videourl}{https://www.youtube.com/watch?v=RcJ1YXHLv5o}
\renewcommand{\videocoverpath}{cover.jpg}

\newcommand{\lecturefigure}[3]{%
\begin{figure}[H]
  \centering
  \includegraphics[width=0.97\textwidth,height=0.49\textheight,keepaspectratio]{slides-images/#1}
  \caption{#2}
  \figsource{#3}
\end{figure}
}

\let\standardtableofcontents\tableofcontents
\renewcommand{\tableofcontents}{%
  \begingroup
  \small
  \setlength{\parskip}{0pt}
  \standardtableofcontents
  \endgroup
}

\begin{document}
\makecscover

\section{Ubicación de la clase: entender Mixtral exige llevar tres cuentas a la vez}

Mixtral suele resumirse en una sola frase: «ocho expertos de 7B, y cada token activa solo dos». La frase basta para atraer la atención, pero no explica por qué el modelo es rápido, por qué ocupa tanta memoria de GPU ni por qué necesita una comunicación compleja, y tampoco aclara qué aprende en realidad cada «experto». El punto de entrada correcto de esta clase es llevar tres cuentas a la vez: el \textbf{capacidad total} registra cuántos parámetros pueden almacenar todos los expertos; el \textbf{cómputo activo} registra cuántos parámetros recorre realmente un token; y el \textbf{costo del sistema} registra la residencia de los pesos, el token dispatch, la comunicación entre tarjetas y el desequilibrio de carga. Ninguna de las tres sustituye a las otras.

\subsection{Primero se corrige la fuente; después se traza la ruta de lectura}

Esta sección aclara primero los límites del material, porque el directorio anterior tenía una contaminación grave de fuentes: el archivo de diapositivas anterior era en realidad el deck de la Lecture 28 de Nathan Lambert sobre alineación, no la clase de Albert Jiang sobre Mixtral. El sitio oficial de CS25 no publicó un deck propio para esta clase, así que aquí se reconstruyen los slide states segundo a segundo a partir del video oficial en 1920x1080, y los 34 minutos de Q\&A se incorporan al texto como teacher voice. Todas las figuras que siguen provienen del video oficial, no del PDF equivocado.

\lecturefigure{slide-01-content.jpg}{Ruta de la clase: dense Transformer, Sparse Mixture of Experts y routing interpretation}{Video oficial de Stanford Online 00:01:59}

\begin{importantbox}{La pregunta central de esta clase}
Sparse Mixture of Experts (SMoE, mezcla dispersa de expertos) no es «un dense model más grande». Amplía el único MLP de cada capa a varios expert MLP, y un router elige un número reducido de expertos para cada token. Así se puede aumentar la capacidad total sin aumentar en la misma proporción los FLOPs por token; sin embargo, todos los expertos siguen teniendo que almacenarse, los tokens siguen teniendo que moverse entre dispositivos y el router todavía puede concentrar la carga en unos pocos expertos.
\end{importantbox}

\begin{knowledgebox}{Términos clave: qué responde cada una de las tres cuentas}
\begin{tabularx}{\textwidth}{L{0.18\textwidth} L{0.28\textwidth} X}
\toprule
Cuenta & Pregunta central & Variables típicas \\
\midrule
Capacidad total & ¿Cuántos pesos almacena el modelo en total? & total parameters, todos los experts, shared attention \\
Cómputo activo & ¿Qué pesos recorre realmente un token? & top-$k$, active parameters, per-token FLOPs \\
Costo del sistema & ¿Cómo se distribuyen los pesos y los tokens en el hardware? & HBM (High Bandwidth Memory, memoria DRAM de alto ancho de banda), expert parallelism, all-to-all, lote, load balance \\
\bottomrule
\end{tabularx}
\end{knowledgebox}

\teachervoice{00:01:04--00:01:54. Albert divide explícitamente la clase en dos partes, architecture e interpretation, e intercala una y otra vez open research questions. Espera que la comunidad abierta siga investigando, en lugar de tratar a Mixtral como una receta ya explicada por completo.}

\subsection{Niveles de evidencia: qué es medición y qué es solo hipótesis de trabajo}

Esta clase reúne configuraciones del modelo, benchmark plots, routing histograms, ablations de la comunidad y juicios basados en la experiencia expresados en el Q\&A. Al leer hay que distinguirlos: la tabla de parámetros del modelo es un hecho estructural; los benchmarks de 2024 son evidencia comparativa de ese momento; «el MLP almacena conocimiento y la attention razona» es la conventional wisdom que el speaker enuncia de forma explícita; que cierto expert se elija con más frecuencia en datos de matemáticas es solo una correlación; y «un expert count mayor facilita la especialización» es un juicio prospectivo. Más adelante se marca el límite en cada caso.

\begin{warningbox}{El expert del nombre no equivale a un especialista profesional humano}
Un expert es solo una subred de parámetros que el router puede seleccionar. Puede ser sensible a la forma de los tokens, a la sintaxis local, a la posición, a la frecuencia o a combinaciones lineales de varias características latentes, y no corresponde de manera natural a un «experto en código» o un «experto en medicina». Traducir directamente la selección del router a semántica humana es el cuarto myth que esta clase busca corregir.
\end{warningbox}

\subsection{Resumen de la sección}

El hilo didáctico de esta clase no es la promoción del modelo, sino reunir sparse capacity, active compute, resident memory, communication e interpretabilidad en un mismo sistema. La corrección de la fuente implica además que toda la slide coverage y las explicaciones orales del docente deben reconstruirse; las notas anteriores no pueden reutilizarse.
\section{Dense baseline: primero, cómo organiza Mistral 7B la attention y la MLP}

Mixtral no reescribe todo el Transformer. Conserva la attention y la residual structure de Mistral 7B y solo sustituye la MLP de cada capa por una expert MLP top-two-of-eight. Para entender con precisión «qué fue lo único que cambió», esta sección establece primero la dense baseline: Grouped-Query Attention, Sliding-Window Attention, RMSNorm, SwiGLU y tensor shape.

\subsection{MHA, MQA y GQA: primero hay que contar los query heads y los KV heads}

La Multi-Head Attention (MHA) estándar da a cada query head su propio key/value head; la Multi-Query Attention (MQA) hace que todos los query heads compartan un solo conjunto de K/V; la Grouped-Query Attention (GQA) se ubica entre ambas: un grupo de query heads comparte un KV head. Mistral 7B usa 32 query heads y 8 KV heads, con el objetivo de conservar más libertad de representación y, al mismo tiempo, reducir la presión sobre la KV cache y la memory bandwidth.

\lecturefigure{slide-02-dense-attention.jpg}{Grouped-Query Attention y Sliding-Window Attention en Mistral 7B}{Grabación oficial de Stanford Online 00:03:55}

Sea $H_q$ el número de query heads, $H_{kv}$ el número de KV heads, $d_h$ la head dimension y $L$ la longitud de la secuencia. En la decodificación autorregresiva, el número de elementos de la KV cache por capa es aproximadamente
\[
N_{KV}=2L H_{kv}d_h,
\]
donde el coeficiente $2$ corresponde a key y value. Si se pasa de $H_{kv}=H_q=32$ en MHA a $H_{kv}=8$ en GQA, con la misma precisión y la misma longitud de secuencia la KV cache se reduce aproximadamente a la cuarta parte.

El núcleo de la attention sigue siendo
\[
\operatorname{Attn}(Q,K,V)=\operatorname{softmax}\!\left(\frac{QK^{\top}}{\sqrt{d_h}}+M\right)V,
\]
donde $M$ es la causal/sliding-window mask. GQA cambia la forma en que se comparten K/V, no esta fórmula de ponderación probabilística.

\begin{knowledgebox}{Concepto previo: por qué la KV cache es un costo de inferencia}
La KV cache guarda los key/value que cada token anterior generó en cada capa, para no recalcular todo el prefijo en cada paso de decodificación. Reside principalmente en la HBM de la GPU (High Bandwidth Memory, la DRAM de video de alto ancho de banda); el crecimiento de la longitud, el número de capas, el número de KV heads o el batch aumenta linealmente su capacidad requerida y consume ancho de banda de lectura y escritura. Por eso el beneficio de GQA no es solo «unas cuantas matrices menos», sino una reducción del estado que permanece residente a largo plazo.
\end{knowledgebox}

\begin{warningbox}{GQA no es una compresión gratuita}
Reducir los KV heads puede hacer perder la independencia de los distintos query heads respecto a la representación de K/V. La calidad final depende de la escala del modelo, los datos de entrenamiento y el head grouping; los resultados de un benchmark no pueden deducirse solo de la fórmula de la caché.
\end{warningbox}

\teachervoice{00:02:18--00:03:36, Albert enfatiza que ni GQA ni la sliding-window attention son invenciones nuevas de Mistral, sino diseños existentes combinados en una dense baseline sólida. Al curso le interesa cómo funcionan en conjunto las decisiones de ingeniería, no adjudicarse una novelty puntual.}

\subsection{Sliding-window attention: cómo se propaga una ventana local a través de las capas}

GQA resuelve el número de KV heads; la Sliding-Window Attention (SWA, atención de ventana deslizante) resuelve el alcance visible de cada capa. Cada posición atiende directamente solo a los $w$ tokens más recientes, de modo que la attention matrix de una capa pasa de $L^2$ conexiones globales a cerca de $Lw$; sin embargo, la información puede propagarse paso a paso a lo largo de las capas, por lo que los tokens de capas profundas pueden absorber indirectamente un contexto más lejano.

El conjunto directamente visible de la posición $t$ en la capa $\ell$ puede escribirse como
\[
\mathcal{N}_{\ell}(t)=\{j\mid \max(0,t-w+1)\le j\le t\}.
\]
Si cada capa permite la transmisión local, la distancia máxima idealizada de propagación de la información crece con el número de capas y es aproximadamente $\ell(w-1)$. Esto no demuestra que «la capa $\ell$ tenga attention global sin pérdidas», pero muestra que el cálculo local puede difundirse mediante la depth.

\begin{importantbox}{GQA y SWA resuelven cuellos de botella distintos}
GQA reduce principalmente la representación de K/V y la KV cache; SWA reduce principalmente las conexiones de attention por capa y el cálculo en secuencias largas. Ambas pueden usarse al mismo tiempo, porque una modifica el head sharing y la otra el token neighborhood.
\end{importantbox}

\subsection{Una capa de dense Transformer: leer las shapes es más confiable que memorizar nombres de módulos}

Los diagramas de arquitectura tienden a ocultar detalles, por eso el ponente presenta directamente una implementación en PyTorch de una sola capa y escribe la tensor dimension en los nombres de las variables. Al leer el código hay que seguir cuatro cosas: la shape de la Q/K/V projection, la codificación posicional después de RoPE, el attention residual y el SwiGLU MLP residual. Solo así se ve que lo único que Mixtral sustituye después es la ruta de la MLP.

\lecturefigure{slide-03-dense-layer-code.jpg}{Tabla de parámetros e implementación en PyTorch de una capa de Transformer de Mistral 7B}{Grabación oficial de Stanford Online 00:06:45}

Si la entrada es $X\in\mathbb{R}^{L\times d}$, el dense block puede resumirse como
\begin{align*}
H &= X+\operatorname{Attention}(\operatorname{RMSNorm}(X)),\\
Y &= H+\operatorname{MLP}(\operatorname{RMSNorm}(H)).
\end{align*}
Mistral usa una MLP de estilo SwiGLU, que puede escribirse como
\[
\operatorname{MLP}(x)=W_2\left(\operatorname{SiLU}(W_1x)\odot W_3x\right),
\]
donde $W_1,W_3$ expanden la hidden dimension a la intermediate dimension, $W_2$ la proyecta de vuelta y $\odot$ es la multiplicación elemento a elemento. Cada expert de Mixtral replica esta parametrización de la MLP, no la capa completa de attention.

\begin{lstlisting}[caption={Pseudocódigo shape-first de un Dense Transformer block}]
q = x @ Wq            # [tokens, query_heads, head_dim]
k = x @ Wk            # [tokens, kv_heads, head_dim]
v = x @ Wv            # [tokens, kv_heads, head_dim]
h = x + attention(q, k, v, causal_window)
y = h + swiglu(rms_norm(h), W1, W2, W3)
\end{lstlisting}

\begin{knowledgebox}{Primer uso: fused kernel}
Un fused kernel (kernel fusionado) combina varias operaciones de GPU, como RMSNorm, projection, activation o gating, en menos kernel launches, lo que reduce las lecturas y escrituras de resultados intermedios en la HBM y el launch overhead. Puede mejorar el rendimiento real, pero no cambia la estructura matemática del modelo; por lo tanto, la optimización de kernels no debe ocultar la diferencia conceptual entre total parameters y active parameters.
\end{knowledgebox}

\teachervoice{00:03:36--00:06:49, el ponente dice que la Transformer architecture tiene una gran cantidad de pequeñas decision choices y que un solo diagrama de bloques no basta; por eso usa código con sufijos de tensor-shape para mostrar a la audiencia qué hace exactamente cada matriz.}

\subsection{Evidencia de la dense baseline: un modelo pequeño no es un modelo débil}

Antes de pasar a los modelos dispersos, el curso usa los benchmarks de Mistral 7B de aquel momento para mostrar que la dense baseline ya era sólida. La comparación de la figura es una instantánea de los modelos al momento de su publicación en 2023, no un leaderboard de 2026; su función es controlar variables: las mejoras de Mixtral deben entenderse en relación con Mistral 7B, de la misma familia, en lugar de atribuir todas las mejoras a MoE.

\lecturefigure{slide-04-mistral7b-performance.jpg}{Comparación de benchmarks de Mistral 7B al momento de su publicación en 2023}{Grabación oficial de Stanford Online 00:07:01}

\begin{knowledgebox}{Lectura de la figura: primero confirmar el tamaño del modelo, luego ver la categoría de la tarea}
La figura compara Mistral 7B con varias configuraciones de Llama 2. El primer paso no es buscar la barra más alta, sino confirmar la escala de parámetros de cada baseline; el segundo es ver si los resultados son consistentes entre tareas; el tercero es revisar si hay ventajas que aparecen en un solo benchmark. El curso solo la usa para demostrar que la dense baseline es competitiva y no toma ninguna gráfica de barras como una clasificación general de capacidades.
\end{knowledgebox}

\begin{warningbox}{Límites de comparabilidad de los benchmarks históricos}
El prompt, la configuración few-shot, el tokenizer, el evaluation harness y la contamination pueden cambiar las puntuaciones. Esta figura sirve para explicar por qué el equipo partió de Mistral 7B para estudiar Mixtral, no para compararlo directamente con modelos más recientes.
\end{warningbox}

\subsection{Resumen de la sección}

Mistral 7B ofrece un dense control claro: GQA reduce los KV heads, SWA limita las conexiones por capa y RMSNorm/SwiGLU organizan el residual block. Mixtral conserva estas estructuras y sustituye la única MLP de la segunda residual branch por una sparse expert mixture. Por lo tanto, todo el análisis posterior de capacidad y costo debe basarse en «attention compartida + varias expert MLP».
\section{Mecanismo de Sparse MoE: cómo el router selecciona y combina los experts}

La sección anterior confirmó que el objeto que se reemplaza es el MLP; en esta sección se aborda el mecanismo de SMoE propiamente dicho. Lo esencial no es copiar ocho veces la red, sino establecer un token-to-expert routing: el router produce un score para cada token, selecciona los top-$k$, normaliza los scores seleccionados y después suma ponderadamente las salidas de varios experts. Todo el proceso debe ser diferenciable, paralelizable y balanceable.

\subsection{MoE no es un concepto nuevo: de la sparsely-gated layer al Switch Transformer}

El curso presenta primero la trayectoria de la literatura y advierte al lector que no confunda el éxito de Mixtral con la primera aparición del concepto de MoE. En 2017, el MoE con compuertas dispersas ya usaba una gating network para seleccionar subredes; el Switch Transformer avanzó en la discusión de data/model/expert parallelism y del entrenamiento disperso a gran escala. Mixtral, por su parte, lleva estas ideas a un modelo decoder-only de alta calidad con pesos abiertos y aporta routing evidence que puede analizarse.

\lecturefigure{slide-05-moe-history.jpg}{Referencias históricas de Mixture of Experts: compuertas dispersas y Switch Transformer}{Video oficial de Stanford Online 00:07:49}

\lecturefigure{slide-06-moe-layer-references.jpg}{Fuentes bibliográficas directas de la Sparse MoE layer}{Video oficial de Stanford Online 00:08:02}

\begin{knowledgebox}{Términos clave: parallelism no es solo «varias GPU»}
\begin{tabularx}{\textwidth}{L{0.19\textwidth} L{0.27\textwidth} X}
\toprule
Modalidad & Objeto que se divide & Comunicación principal \\
\midrule
Data parallelism & Distintos shards del lote; el modelo se replica & gradient all-reduce \\
Tensor/model parallelism & Matrices de una sola capa o la hidden dimension & all-reduce / all-gather \\
Expert parallelism & Distintos experts colocados en distintos ranks & token all-to-all dispatch / combine \\
Pipeline parallelism & Distintas capas colocadas en distintos stages & activation send/receive \\
\bottomrule
\end{tabularx}
Las collectives son primitivas de comunicación colectiva entre varias GPU, por ejemplo all-reduce, all-gather, reduce-scatter y all-to-all. La comunicación más representativa de MoE es all-to-all: cada GPU envía los tokens a la GPU que aloja el expert de destino y, tras el cálculo, devuelve los resultados.
\end{knowledgebox}

\teachervoice{00:07:01--00:08:24, Albert repite que MoE es «not new». Recomienda en especial la discusión sobre parallelism del Switch Transformer y sitúa a Mixtral dentro de la genealogía continua de la sparsely-gated layer de 2017.}

\subsection{Top-two routing: de los logits a la weighted expert output}

Para una token representation $x\in\mathbb{R}^{d}$, el router usa la matriz $W_r\in\mathbb{R}^{E\times d}$ para producir $E$ expert logits:
\[
z=W_rx.
\]
Sea $S_k(x)$ el conjunto de los $k$ expert index con los valores más altos de $z$; Mixtral usa $E=8,k=2$. La normalización se hace solo sobre el conjunto seleccionado:
\[
g_i(x)=
\begin{cases}
\dfrac{\exp(z_i)}{\sum_{j\in S_k(x)}\exp(z_j)}, & i\in S_k(x),\\[0.8em]
0, & \text{otherwise}.
\end{cases}
\]
La salida final del MLP de MoE es
\[
y=\sum_{i=1}^{E}g_i(x)\,\operatorname{Expert}_i(x).
\]
Estos tres pasos responden, respectivamente, a «puntuar», «seleccionar» y «combinar». Sparse significa que solo se ejecutan $k$ experts, no que los pesos restantes desaparezcan del modelo.

\lecturefigure{slide-07-moe-layer-diagram.jpg}{El router selecciona experts para cada token y combina sus salidas}{Video oficial de Stanford Online 00:08:19}

\begin{importantbox}{Semántica algorítmica del top-$k$ routing}
Cada token se enruta de forma independiente; tokens adyacentes de una misma secuencia pueden seleccionar experts distintos; un mismo expert recibe simultáneamente tokens provenientes de muchas secuencias. La implementación debe reordenar los tokens por expert, hacer el dispatch entre GPU, ejecutar por lotes los expert MLP y después aplicar un inverse-permute para restaurar el orden original.
\end{importantbox}

\begin{lstlisting}[caption={Pseudocódigo mínimo del top-two token routing}]
router_logits = hidden @ router_weight.T
expert_ids = topk(router_logits, k=2)
gates = softmax(gather(router_logits, expert_ids), dim=-1)
expert_inputs = dispatch_by_expert(hidden, expert_ids)
expert_outputs = run_selected_experts(expert_inputs)
output = combine_and_restore_order(expert_outputs, gates)
\end{lstlisting}

\teachervoice{00:08:24--00:09:18, el ponente lo desglosa oralmente en el mismo orden: el input entra primero al router, el router selecciona los top two y asigna gating weights, cada expert procesa de forma independiente y luego se hace la suma ponderada con esos weights.}

\subsection{Página de fórmulas y tabla de parámetros: cómo leer \texorpdfstring{$E=8$, $k=2$}{E=8, k=2} y las shared parts}

La página de fórmulas reúne en una misma slide el mecanismo y la configuración del modelo. Al leer la tabla hay que distinguir \texttt{num\_experts=8} de \texttt{top\_k\_experts=2}: el primero determina la capacidad disponible en cada capa; el segundo, la ruta activa de un token. Attention, embedding, router y residual structure son shared parameters y no se multiplican simplemente por ocho según el número de experts.

\lecturefigure{slide-08-moe-equation-table.jpg}{Fórmula del top-two routing de Mixtral y tabla de parámetros del modelo}{Video oficial de Stanford Online 00:09:27}

Si se denota con $P_s$ el shared parameter de cada capa, con $P_e$ los parámetros de un solo expert MLP y con $N$ el número de capas, los parámetros totales y activos aproximados son
\begin{align*}
P_{\text{total}} &\approx P_{\text{embed}}+N(P_s+E P_e),\\
P_{\text{active}} &\approx P_{\text{embed}}+N(P_s+k P_e).
\end{align*}
Por lo tanto, el $8\times 7\mathrm{B}$ del nombre del modelo no es ni el total exacto de parámetros ni los parámetros activados por token; las shared parts hacen que la simple multiplicación no sea exacta.

\begin{warningbox}{Lo «disperso» ocurre en la ruta de ejecución, no en la existencia en almacenamiento}
Los experts no seleccionados no se ejecutan para el token actual, pero siguen formando parte del checkpoint y, por lo general, deben permanecer residentes en la HBM de la GPU, la DRAM de la CPU u otro nivel de memoria. La resident memory la determinan los total parameters; los MLP FLOPs por token se acercan más a lo que determinan los active parameters.
\end{warningbox}

\subsection{Mixtral 8x7B: la convergencia formal de capacidad y ruta activa}

Ahora puede leerse la slide completa del modelo. Mixtral hereda el attention block de Mistral 7B, cada capa tiene ocho SwiGLU experts y cada token selecciona dos. Las cifras release-time que da el curso son de aproximadamente 46.7B total parameters y 12.9B active parameters; corresponden, respectivamente, a la capacidad y residencia en memoria y a la ruta de ejecución por token, y no son intercambiables.

\lecturefigure{slide-09-mixtral-8x7b.jpg}{Estructura del modelo Mixtral 8x7B, tabla de parámetros y resumen de capacidades al momento del lanzamiento}{Video oficial de Stanford Online 00:10:26}

\begin{importantbox}{Una descripción precisa de Mixtral}
Mixtral 8x7B es un SMoE decoder-only con «attention y router compartidos, ocho expert MLP por capa y activación top-two por token». No es una votación entre ocho modelos 7B independientes, ni un dense model que solo ocupe memoria para 13B de pesos.
\end{importantbox}

\teachervoice{00:09:38--00:10:17, Albert usa la active path para explicar la cost-performance frontier y enumera además las características de entonces: soporte multilingüe, 32K context, Apache 2.0 y los benchmark claim. Estas notas las conservan como evidencia release-time de 2024, no como especificaciones actuales del producto.}

\subsection{Resumen de la sección}

La estructura matemática de Sparse MoE es breve: router logits, top-$k$ mask, selected softmax y expert weighted sum; lo verdaderamente difícil es llevar este algoritmo a nivel de token a varias GPU y mantener el balance de carga. La descomposición entre parámetros totales y activos es la base común de todas las discusiones posteriores sobre myth, rendimiento y despliegue.
\section{Por qué MoE-fy MLP: evidencia de desempeño e hipótesis de trabajo}

La sección anterior explicó «cómo enrutar»; esta pregunta «por qué conviene ampliar primero el MLP». El ponente adopta una intuición común, aunque no demostrada, sobre la división de funciones: el MLP se parece más a un almacén de conocimiento y la attention, a un mecanismo de composición y de algoritmos. Si esta intuición es parcialmente cierta, ampliar la MLP capacity debería ayudar más en las knowledge-heavy tasks; a continuación, la clase pone a prueba esta tendencia con dos grupos de gráficas de benchmark y plantea la pregunta inversa de MoE-fying attention.

\subsection{El MLP almacena conocimiento y la attention ejecuta algoritmos: una intuición útil pero riesgosa}

El valor de esta slide está en que enuncia la hipótesis de forma explícita, no en que la demuestre. El conocimiento y el cómputo de un Transformer se distribuyen entre embeddings, attention, MLP, normalization y residual interactions; dividir las funciones en dos categorías sirve para organizar experimentos, pero no permite afirmar que una layer o un expert se encarga de una sola función cognitiva.

\lecturefigure{slide-10-why-moe-mlp.jpg}{Por qué conviene MoE-fy primero el MLP: la hipótesis de trabajo knowledge-versus-reasoning}{Grabación oficial de Stanford Online 00:10:58}

\begin{warningbox}{No conviertas una intuición sobre el mecanismo en un teorema de localización}
Aunque la ampliación del MLP y la mejora en el knowledge benchmark ocurran al mismo tiempo, no puede concluirse que «el conocimiento solo existe en el MLP». El aumento de parámetros, los cambios en la dinámica de entrenamiento, la router regularization y el aprovechamiento de los datos pueden influir en el resultado. La formulación confiable es: esta clase usa esa intuición para formular una predicción y luego examina si las categorías de tareas la cumplen aproximadamente.
\end{warningbox}

\teachervoice{00:10:32--00:11:44, Albert llama directamente conventional wisdom a esta distinción y anticipa que MoE-fying MLP mejorará sobre todo la knowledge; no la presenta como una conclusión rigurosa de neurociencia.}

\subsection{Category bars: la mayor ganancia está en las tareas de conocimiento, pero no es la única}

La gráfica de barras compara Mistral 7B, Mixtral 8x7B y varias baselines de Llama 2. Al leerla, primero agrupa por category y después compara la diferencia dense/sparse dentro de la misma familia; no empieces por el promedio general. El ponente observa que la mejora más notable está en las knowledge-heavy tasks; reasoning/comprehension también mejoran, aunque en menor medida.

\lecturefigure{slide-11-mixtral-performance-bars.jpg}{Benchmark de Mixtral 8x7B al momento de su release-time en distintas categorías de tareas}{Grabación oficial de Stanford Online 00:11:41}

Si la puntuación en una tarea es $s_m$ y la de la baseline correspondiente es $s_b$, puede usarse la ganancia relativa
\[
\Delta_{\mathrm{rel}}=\frac{s_m-s_b}{|s_b|+\epsilon}
\]
para comparar magnitudes distintas, aunque esto sigue sin resolver la benchmark saturation, la prompt sensitivity ni la contamination. La conclusión didáctica de la figura es que «la ganancia se distribuye de manera desigual», no que «el conocimiento y el razonamiento ya están completamente separados».

\begin{knowledgebox}{Lectura de la figura: cuatro pasos para evitar el category cherry-picking}
Primero, confirma si cada grupo usa el mismo evaluation harness; segundo, distingue métricas como accuracy y exact match; tercero, observa la diferencia de Mistral 7B a Mixtral dentro de la misma familia; cuarto, revisa si el patrón es consistente entre categorías. Elegir solo la barra más alta exagera la generalidad de MoE.
\end{knowledgebox}

\subsection{Active-parameter plot: el eje horizontal no es el total model size}

El segundo grupo de gráficas pone active parameters en el eje horizontal y subraya que Mixtral usa unos 12.9B parámetros por token y, aun así, se ubica en una mejor performance region. Este eje responde «cuánta calidad se obtiene con un cómputo activo similar», pero no responde «qué tan grande es el checkpoint», «cuánta memoria de GPU ocupa» ni «cuánto cuesta el all-to-all». Por eso respalda el argumento capacity-to-compute, pero no el argumento capacity-to-memory.

\lecturefigure{slide-12-mixtral-performance-plots.jpg}{Gráfica cost-performance de active parameters contra el desempeño en distintas tareas}{Grabación oficial de Stanford Online 00:12:34}

Un serving cost más completo, aunque todavía simplificado, puede escribirse como
\[
C_{\text{serve}}\approx C_{\text{shared}}+kC_{\text{expert}}+C_{\text{router}}+C_{\text{dispatch}}+C_{\text{imbalance}},
\]
donde los dos primeros términos corresponden aproximadamente al cómputo y los tres últimos, al enrutamiento, la comunicación y la espera. Los active parameters solo aproximan los dos primeros términos; no representan toda la latency ni el dollar cost.

\begin{warningbox}{La Pareto frontier de la figura no incluye todas las variables de producción}
La gráfica de Pareto formada por benchmark quality y active parameters no considera HBM residency, interconnect, batch size, quantization, kernel maturity ni tail latency. La selección de ingeniería requiere una contabilidad de sistema aparte.
\end{warningbox}

\teachervoice{00:11:44--00:12:38, el ponente explica personalmente que el eje horizontal son los active parameters y el vertical, el desempeño en cada tipo de tarea; subraya que la mejora en la categoría knowledge de Mistral 7B a Mixtral es especialmente grande.}

\subsection{Por qué no MoE-fy attention: la estabilidad antes que la imaginación}

Si ampliar la MLP capacity funciona, la pregunta natural es convertir también las Q/K/V projection en switch layers. La clase señala que esta línea ya se ha explorado experimentalmente, pero la estabilidad numérica es un obstáculo real: algunas configuraciones se pueden entrenar en fp32 y divergen (diverge) al pasar a bf16. Como el entrenamiento de los grandes modelos actuales depende de la eficiencia de la baja precisión, poder entrenar de forma estable importa más que el ahorro de parámetros sobre el papel.

\lecturefigure{slide-13-moe-attention-question.jpg}{Pregunta abierta: convertir también la attention projection en MoE}{Grabación oficial de Stanford Online 00:14:06}

\begin{knowledgebox}{Concepto previo: bf16 stability}
bf16 conserva el mismo número de exponent bits que fp32, pero su mantissa es más corta, por lo que los errores de redondeo en multiplicaciones y sumas, softmax, normalization y router logits son más notorios. Que el entrenamiento diverja (diverge) depende de scaling, normalization, optimizer, initialization y kernel implementation; no puede atribuirse solo a que «la baja precisión es mala».
\end{knowledgebox}

\teachervoice{00:12:38--00:14:03, Albert deja MoE attention como una open research question explícita: los intentos existentes pueden divergir (diverge) en bf16, y en el futuro se necesitarán mejores técnicas de normalization o de estabilización.}

\subsection{Resumen de la sección}

Las gráficas de desempeño respaldan una conclusión limitada: con active parameters similares, el capacity-to-compute tradeoff de Mixtral es muy favorable, y la ganancia en las knowledge-heavy tasks destaca especialmente. No demuestran que el MLP concentre en exclusiva el conocimiento, ni que los active parameters equivalgan al costo real de servicio. MoE-fying attention, por su parte, nos recuerda que toda ampliación estructural debe superar el umbral de ingeniería de la estabilidad en baja precisión.
\section{Cuatro myths: del nombre del modelo a la contabilidad del sistema}

La parte de mayor valor didáctico de la clase es el desmontaje sucesivo de cuatro errores intuitivos. Cada uno confunde, respectivamente, expert identity, total parameters, active compute y semantic specialization. Esta sección ofrece para cada uno una explicación alternativa que puede calcularse, e incorpora lo que se dijo en la sesión de Q\&A sobre memory, communication y throughput.

\subsection{Myth 1: Mixtral solo tiene ocho experts globales}

La sección anterior ya explicó que Mixtral introduce experts dispersos únicamente en la ruta MLP de cada capa; aquí se corrige primero la imagen de roles globales que sugiere el nombre «8x7B». Al leer el diagrama de estructura siguiente, conviene distinguir primero entre «cuántos MLP seleccionables hay en cada capa» y «cuántas identidades rastreables entre capas tiene el modelo completo», y solo después juzgar si el número de un expert puede tener contenido semántico. «8x7B» invita a imaginar ocho expertos que atraviesan todas las capas, pero en realidad cada Transformer layer tiene ocho expert MLP; el expert 0 de capas distintas no comparte identity. Con $N=32$ capas, la estructura contiene $32\times 8=256$ layer-local expert modules.

\lecturefigure{slide-14-myth-eight-experts.jpg}{Myth 1: cada capa tiene ocho experts, no ocho experts globales}{Grabación oficial de Stanford Online 00:15:12}

Dentro de una misma capa, las expert labels tienen permutation symmetry. Para cualquier expert permutation $\pi$, si se permutan a la vez los expert weights y los output channels del router, la función no cambia:
\[
\sum_i g_i(x)E_i(x)=\sum_i g_{\pi(i)}'(x)E_{\pi(i)}'(x).
\]
Por lo tanto, «expert 3» solo tiene sentido con un checkpoint, una layer y un labeling fijos; no se pueden comparar los números entre capas ni entre runs de entrenamiento.

\teachervoice{00:14:15--00:15:09, Albert señala que el nombre del modelo puede inducir a error: hay ocho expertos por capa y su numeración dentro de la capa es permutable; en total hay muchos layer-local experts, no ocho roles permanentes.}

\begin{warningbox}{No rastrees el número de un expert entre capas}
El expert 3 de la layer 15 y el expert 3 de la layer 31 no tienen continuidad semántica intrínseca. Para estudiar «la evolución de un mismo expert», primero hay que definir un alignment comparable, en lugar de comparar solo índices de arreglo.
\end{warningbox}

\subsection{Myth 2: \texorpdfstring{$8\times 7\mathrm{B}=56\mathrm{B}$}{8 x 7B = 56B} es el número exacto de parámetros}

La multiplicación directa ignora los shared embeddings, la attention, la normalization y el router. La clase da aproximadamente 46.7B total parameters y 12.9B active parameters. Lo que realmente hay que recordar no son las cifras decimales, sino la decomposition: los parámetros totales cuentan todos los expert weights; los parámetros activos cuentan solo los dos experts que eligió el token actual, más las shared parts que siempre se ejecutan.

\lecturefigure{slide-15-myth-56b.jpg}{Myth 2: la shared attention hace que el total de parámetros no sea simplemente 56B}{Grabación oficial de Stanford Online 00:15:37}

Se define la expert fraction
\[
\rho=\frac{kP_e}{P_s+kP_e},
\]
que expresa la proporción de la ruta activa que corresponde a los expert MLP. Aunque $k/E=1/4$, la razón active/total global no es exactamente un cuarto, porque los shared parameters aparecen tanto en el numerador como en el denominador.

\begin{importantbox}{Formato mínimo completo para reportar el número de parámetros}
Hay que reportar al menos, de forma conjunta: total parameters, active parameters per token, number of experts $E$, top-$k$, la configuración de la shared attention y la precisión. Escribir solo «13B-class» oculta que 46.7B de pesos deben permanecer en memoria; escribir solo «47B model» oculta la ejecución dispersa por token.
\end{importantbox}

\teachervoice{00:15:09--00:15:41, el ponente corrige la cifra de 56B: la attention y el gating son compartidos, y el total es de aproximadamente 46.7B; cada token ve unos 12.9B active parameters, no 14B.}

\subsection{Myth 3: el cost es proporcional a los active parameters}

Si el modelo se ejecuta secuencialmente en una sola tarjeta, los active parameters son un proxy importante de los FLOPs; pero en cuanto los experts se reparten mediante sharding entre varias tarjetas, los tokens deben reordenarse y enviarse según la decisión del router. Sharding (particionamiento) significa colocar los expert weights en distintos devices, de modo que cada tarjeta ya no replique todos los experts; ahorra almacenamiento por tarjeta, pero introduce dispatch/combine communication.

\lecturefigure{slide-16-myth-cost-active.jpg}{Myth 3: los active parameters no representan todo el serving cost}{Grabación oficial de Stanford Online 00:16:42}

Si un lote tiene $T$ tokens y el expert $i$ recibe $n_i$ tokens, la carga uniforme ideal es $T k/E$. Se define el imbalance ratio
\[
I=\frac{\max_i n_i}{Tk/E}.
\]
$I=1$ indica la distribución más uniforme; $I>1$ indica que el expert más ocupado supera el promedio. El step paralelo suele quedar determinado por $\max_i n_i$, de modo que unos pocos puntos calientes pueden dejar a los demás devices esperando sin trabajo.

\begin{knowledgebox}{Primer uso: expert parallel all-to-all}
La expert parallelism reparte los experts entre ranks. El dispatch all-to-all envía primero cada token al rank donde está su expert de destino; tras el cálculo del expert, el combine all-to-all devuelve la salida al rank original. La comunicación entre nodes suele ser más lenta que el GPU interconnect dentro de un mismo node, por lo que la topology y el placement modifican la latency.
\end{knowledgebox}

\teachervoice{00:15:41--00:16:49, Albert rechaza explícitamente vincular el cost con el active parameter count. El routing dinámico exige enviar tokens y no permite fijar las rutas de antemano, por lo que el serving real tiene un costo de comunicación adicional respecto de un dense model con el mismo active-size.}

\subsection{Inference load balance: el expert más lento determina cuándo termina la ronda}

El balanceo de carga no es solo una training auxiliary loss. En inference, la prompt distribution real puede inclinar al router hacia unos pocos experts y producir un straggler. Aunque la cantidad promedio de cálculo sea correcta, basta con que la queue de un expert sea demasiado larga para que la fase de combine de todo el lote tenga que esperar a que termine.

\lecturefigure{slide-17-inference-load-balance.jpg}{Problema abierto: cómo balancear las expert loads en inference}{Grabación oficial de Stanford Online 00:17:36}

La capacity suele expresarse como
\[
C=\left\lceil \alpha\frac{Tk}{E}\right\rceil,
\]
donde $\alpha\ge 1$ es el capacity factor. Si un expert recibe más de $C$ tokens, puede hacer drop, reroute o diferir su procesamiento; cada estrategia implica un intercambio entre quality, latency y determinism.

\begin{warningbox}{Una distribución uniforme en entrenamiento no garantiza una distribución uniforme en producción}
La router distribution sobre la training data puede ser casi uniforme y, aun así, los prompts de producción pueden concentrarse en ciertos patrones. En producción conviene registrar el per-expert token count, el queue time, la drop/reroute rate y la tail latency, en lugar de mirar solo el promedio de tokens/expert.
\end{warningbox}

\teachervoice{00:16:49--00:17:36, el ponente describe la inference load balance como un problema de esperar al expert más lento y menciona exploraciones como capacity-aware neighbor routing y mixture of depths.}

\subsection{Compressing SMoEs: cuantización, offload y sparsification no son lo mismo}

La clase presenta la compresión como una open question, no como una recipe ya resuelta. La quantization (cuantización) reduce el número de bits de cada weight; el CPU offload traslada parte de los experts de la HBM de la GPU a la DRAM de la CPU y los transfiere cuando se necesitan; la sparsification intenta eliminar o poner en cero los pesos poco importantes dentro de cada expert. Las tres optimizan recursos distintos y tienen costos distintos.

\lecturefigure{slide-18-compressing-smoe.jpg}{Problema abierto: ¿se puede comprimir una SMoE a una memoria mínima?}{Grabación oficial de Stanford Online 00:19:07}

Si los total parameters son $P$ y cada parámetro ocupa $b$ bits, la cota inferior ideal del almacenamiento de pesos es
\[
M_{\text{weights}}=\frac{Pb}{8}\ \text{bytes}.
\]
En la práctica hay que sumar scale/zero-point, KV cache, activation workspace, router buffers y framework metadata. El offload, por su parte, añade un tiempo de transferencia aproximado
\[
t_{\text{transfer}}\approx \frac{M_{\text{moved}}}{B_{\text{CPU--GPU}}},
\]
donde $B_{\text{CPU--GPU}}$ es el ancho de banda efectivo de la interconexión.

\begin{knowledgebox}{Términos clave: tres tipos de compresión}
\begin{tabularx}{\textwidth}{L{0.18\textwidth} L{0.24\textwidth} X}
\toprule
Método & Ahorro principal & Riesgo principal \\
\midrule
Quantization & bytes por weight & pérdida de precisión, dequant kernel, scale overhead \\
Offload & GPU resident memory & PCIe/NVLink transfer, cache miss, tail latency \\
Sparsification & nonzero weights/FLOPs efectivos & kernels irregulares, degradación de calidad, reentrenamiento \\
\bottomrule
\end{tabularx}
\end{knowledgebox}

\teachervoice{00:17:36--00:19:15, Albert cita conjeturas optimistas de la comunidad sobre la compresión extrema, pero reconoce que, al momento de la clase, aún no había visto resultados convincentes de extreme SMoE sparsification; se trata de una agenda de investigación, no de un compromiso de lanzamiento.}

\subsection{Resumen de la sección}

Los cuatro myths comparten un mismo origen: tomar el nombre por el modelo del sistema. La alternativa correcta es la siguiente: los experts son layer-local y permutables; los total parameters y los active parameters deben separarse; el serving cost incluye además communication e imbalance; la semantic domain specialization no es necesaria ni queda demostrada por los diagramas de routing. El load balance y la compression, por su parte, llevan estos conceptos a las restricciones reales del hardware.
\section{Routing interpretation: el gate discreto aporta evidencia, pero no proporciona semántica por sí mismo}

Las activation de una dense network se ubican en un espacio continuo de alta dimensión y son difíciles de interpretar directamente; el SMoE router, al menos, entrega expert IDs discretos y gating weights, lo que permite a los investigadores contabilizar la asignación de tokens a experts. Esta sección sigue los tres niveles de análisis del curso: domain histogram, consecutive-token persistence y token-level map; al final se discuten la expert ablation hecha por la comunidad y la conclusión de que «un expert no es un domain».

\subsection{Por qué el discrete gate parece más interpretable}

En cada capa, el router genera para cada token un conjunto pequeño $S_k(x)$, que resulta más fácil de contar que una activation de miles de dimensiones. Los investigadores pueden preguntarse: ¿los distintos domain de datos tienden hacia experts distintos?, ¿los tokens adyacentes conservan el mismo ruteo?, ¿qué ocurre si se elimina cierto expert? Sin embargo, estas preguntas solo producen observable correlations; la interpretación todavía requiere experimentos de control.

\lecturefigure{slide-19-discrete-gating.jpg}{Las discrete gating signals aportan nuevas magnitudes observables para la SMoE interpretation}{Grabación oficial de Stanford Online 00:20:19}

Para el domain $d$, la layer $\ell$ y el expert $i$, la routing frequency puede escribirse como
\[
p_{\ell}(i\mid d)=\frac{\#\{x\in d:i\in S_k^{(\ell)}(x)\}}{k\,|d|}.
\]
Si fuera completamente uniforme, $p_{\ell}(i\mid d)\approx 1/E$. Una desviación respecto de la uniformidad solo significa que «los tokens de ese domain eligen con más frecuencia a ese expert», no que «la semántica de ese expert sea ese domain».

\begin{warningbox}{Ni la attention weight ni el router ID son explicaciones causales}
La visualización nos dice qué miró el modelo o a quién eligió, pero eso no equivale a que dicho componente tenga un papel causal único sobre la salida. Se necesitan intervention, ablation, counterfactual routing o retraining para acercarse a una conclusión causal.
\end{warningbox}

\teachervoice{00:19:15--00:20:24, el ponente considera que el discrete gate ofrece, frente a la dense activation, una «incredible opportunity», pero su pregunta sigue siendo «can we make some sense out of this», y no una declaración de que la interpretabilidad ya esté resuelta.}

\subsection{Domain histogram: el spike de las capas intermedias es interesante, pero la evidencia sigue siendo débil}

El experimento introduce en el modelo distintos validation subsets de The Pile y contabiliza la expert selection en las layer 0, 15 y 31. Las capas superficiales y profundas son, en general, cercanas a la uniformidad; en una capa intermedia, cierto expert presenta una tasa de selección más alta para DeepMind Mathematics/GitHub. Al leer la figura, primero se observa la baseline $1/E=12.5\%$ y después se comparan los cambios entre capas, en lugar de ponerles nombre a los colores directamente.

\lecturefigure{slide-20-domain-specialisation.jpg}{Distribución de expert routing de los distintos domains de The Pile en las capas superficiales, intermedias y profundas}{Grabación oficial de Stanford Online 00:22:39}

\begin{knowledgebox}{Lectura de la figura: qué podría verse en cada una de las tres capas}
La layer 0 está cerca del raw token, por lo que el ruteo puede verse influido por caracteres, signos de puntuación y sintaxis local; las capas intermedias contienen más representaciones combinadas, de ahí que aparezca un spike correlacionado con el domain; la última capa está cerca del decoding objective, y su distribución puede volver a tender hacia necesidades de salida compartidas. Esta narración es una hypothesis comprobable, no algo que una sola figura determine de manera única.
\end{knowledgebox}

\teachervoice{00:20:24--00:22:28, Albert dice que el spike de math/code del expert 3 en la layer 15 es «interesting», y enseguida añade que es highly speculative y que cannot conclude much. Esta autolimitación debe conservarse.}

\subsection{Consecutive tokens: la routing persistence es más fuerte en las capas intermedias}

El domain histogram solo indica la proporción global de selección en distintos subconjuntos de datos; no permite saber si el router mantiene su decisión dentro de un fragmento local. Por eso, el segundo experimento cambia la pregunta de «qué domain elige con más frecuencia a quién» a «si los tokens adyacentes eligen de forma consecutiva al mismo expert». Al leer la figura, primero se confirma la baseline aleatoria de los dos eventos, top-1 y top-2, y después se comparan las capas superficiales, intermedias y profundas, sin interpretar directamente la continuidad como especialización temática. Si el top-1 fuera independiente y uniforme entre ocho experts, la baseline aleatoria sería $1/8=12.5\%$; si se compara «si el siguiente token coloca a ese expert en su top-2», la baseline es más alta. En la figura, la persistence de la layer 15 es claramente superior a la aleatoria, la de la layer 0 es más débil y la de la layer 31 vuelve a bajar.

\lecturefigure{slide-21-consecutive-tokens.jpg}{Frecuencia con la que tokens adyacentes siguen eligiendo al mismo expert}{Grabación oficial de Stanford Online 00:25:13}

Se define la top-1 persistence
\[
P_{\ell}^{(1)}=\Pr\!\left(e_{t+1}^{(1)}=e_t^{(1)}\right),
\]
así como la set persistence
\[
P_{\ell}^{(2)}=\Pr\!\left(e_t^{(1)}\in S_2(x_{t+1})\right).
\]
Ambas miden la continuidad local del ruteo, pero no demuestran que los tokens consecutivos compartan un «experto temático»; subword, syntax, position o hidden-state smoothness pueden causar persistence.

\teachervoice{00:22:46--00:25:18, el ponente explica capa por capa la random baseline y la observed frequency: la layer 15 es aproximadamente la más fuerte, la última capa sigue por encima de lo aleatorio pero desciende, y afirma explícitamente que se requiere un análisis más fino.}

\subsection{Token map: los colores parecen mostrar patrones, pero la semántica no es ordenada}

El tercer experimento colorea por expert los tokens de código, aritmética y preguntas de opción múltiple. Algunos números o fragmentos locales aparecen con colores similares, pero en conjunto no hay una división clara del tipo «todo el código va al expert A y todas las matemáticas al expert B». La función didáctica más importante de la imagen es que el lector vea por sí mismo la fragmentación del expert routing, no construir una taxonomy atractiva.

\lecturefigure{slide-22-token-assignment.jpg}{Token-level expert assignment en muestras de código, aritmética y preguntas y respuestas}{Grabación oficial de Stanford Online 00:26:13}

\begin{knowledgebox}{Lectura de la figura: primero la continuidad local, luego la consistencia entre muestras}
Primer paso: buscar bloques de color consecutivos dentro de una misma secuencia; segundo paso: comparar si tokens del mismo tipo conservan el color en muestras distintas; tercer paso: comparar entre layers. Solo si el patrón es estable en múltiples muestras, entre capas y entre semillas aleatorias vale la pena intentar asignar una semántica a un expert.
\end{knowledgebox}

\teachervoice{00:25:18--00:26:19, Albert observa que algunos digits se envían al mismo expert, pero que en general «doesn't seem to be much specialization». El curso no deduce del token map ninguna domain taxonomy clara.}

\subsection{Myth 4: por qué los domain experts puros también podrían reducir la tasa de utilización}

Supongamos que solo dos experts se especializan en code: entonces una secuencia larga de código concentraría gran cantidad de tokens en esos dos experts, mientras los otros seis quedan ociosos; esto produce load imbalance y además desperdicia capacity paralela. Una specialization más razonable podría consistir en latent features compartidas entre domains, de modo que todos los experts puedan aprovecharse con distintas entradas.

\lecturefigure{slide-23-myth-domain-specialisation.jpg}{Myth 4: los experts óptimos no necesariamente se reparten el trabajo según los domains humanos}{Grabación oficial de Stanford Online 00:27:15}

Si todos los tokens del domain $d$ se rutean al subconjunto $A_d$, la cota superior de la tasa de utilización es aproximadamente
\[
U_d=\frac{|A_d|}{E}.
\]
Por ejemplo, con $|A_d|=2,E=8$, un batch de un solo domain aprovecha plenamente, como máximo, una cuarta parte de los experts. Un router real puede combinar top-two y repartir el trabajo entre features, pero la fórmula revela la tensión entre «expertos de domain limpios» y una alta tasa de utilización.

\begin{importantbox}{Una definición más prudente de specialization}
No hay que asignar primero un nombre a un expert y luego buscar ejemplos; hay que medir la distribución del ruteo, el representation probe, la causal ablation y el load behavior, y después describir a qué features es sensible cierto expert. Una feature puede aparecer simultáneamente en código, matemáticas y lenguaje natural.
\end{importantbox}

\teachervoice{00:26:20--00:27:22, el ponente usa el contraejemplo de los coding experts para explicar que, si los tokens de código solo fueran a dos experts, los demás experts se quedarían parados sin trabajar; el lenguaje es demasiado complejo como para reducir la specialization a domains de alto nivel.}

\subsection{Treasure hunt: los pesos abiertos permiten que la comunidad haga intervention con rapidez}

Aproximadamente un día después de la publicación del modelo, la comunidad intentó retirar experts específicos y descubrió que eliminar cierto número hacía que el MMLU se desplomara drásticamente, lo que se convirtió en un meme de «uno trabaja y los demás miran». El resultado es llamativo, pero podría reflejar la router calibration, un layer-specific bottleneck, la implementation o el método de distributed ablation; no prueba directamente que un solo expert almacene todo el conocimiento.

\lecturefigure{slide-24-treasure-hunt.jpg}{Expert ablation de la comunidad, desplome del MMLU y meme del código abierto}{Grabación oficial de Stanford Online 00:28:31}

Se define el metric change tras eliminar el expert $i$
\[
\Delta_i=m(\theta_{-i})-m(\theta),
\]
donde $\theta_{-i}$ denota el modelo intervenido con cierto expert deshabilitado. Un $\Delta_i$ negativo y grande demuestra que la intervención afecta la salida, pero no permite distinguir entre el conocimiento del propio expert, la router re-normalization o el downstream distribution shift.

\begin{warningbox}{El diseño de una ablation debe registrar cuatro cosas}
¿Se eliminó el expert de una sola layer o el del mismo número en todas las capas? ¿Se renormalizaron los Gates? ¿Cómo se manejan los overflow token? ¿La Metric es sensible al formato de salida? Si no se registra esto, los resultados de la comunidad solo pueden ser pistas.
\end{warningbox}

\teachervoice{00:27:22--00:28:38, Albert toma este hallazgo de la comunidad, surgido en menos de 24 horas, como ejemplo de «why we love open source»; su valor está en la intervention rápida, no en que el meme ya explique el modelo.}

\subsection{Agenda final de interpretación: buscar latent feature subspaces}

Al final, el curso no les pone etiquetas a los experts, sino que plantea una pregunta más difícil: los experts podrían codificar combinaciones lineales de conceptos humanos o ser sensibles a cierto feature subspace. La investigación futura necesita combinar las routing stats con representation analysis, causal patching, expert swapping y controlled retraining.

\lecturefigure{slide-25-interpret-moe-decisions.jpg}{Pregunta abierta: qué features aprenden realmente los experts}{Grabación oficial de Stanford Online 00:29:11}

\begin{importantbox}{La escalera de evidencia de la visualización a la explicación}
Routing histogram $\rightarrow$ token map $\rightarrow$ controlled ablation $\rightarrow$ counterfactual rerouting $\rightarrow$ retraining/feature recovery. Cuanto más se avanza, más cerca se está de la causalidad, pero también mayor es el costo. Un solo histogram solo puede ubicarse en el punto de partida de la escalera.
\end{importantbox}

\teachervoice{00:28:38--00:29:30, el ponente considera que los experts podrían capturar features distintas de los conceptos humanos, incluso combinaciones lineales de varios conceptos; la verdadera pregunta es cómo recuperar su subspace.}

\subsection{Resumen de la sección}

El Discrete routing efectivamente añade variables observables a la interpretabilidad: domain frequency, routing persistence, token assignment y ablation sensitivity. Pero todos los resultados apuntan a la misma conclusión prudente: los experts presentan un comportamiento estructurado, pero no se reparten el trabajo según domains humanos simples. La dirección de investigación más valiosa es combinar el gate discreto con la intervención causal, la representation geometry y la carga del sistema.
\section{Síntesis de ingeniería de la sesión de Q\&A: cuándo vale la pena MoE y cuándo dense es más honesto}

La sesión de Q\&A posterior a la exposición principal ocupa más de la mitad de la grabación y aporta las condiciones de despliegue que las slides no desarrollan. Esta sección no acumula las respuestas en el orden de las preguntas, sino que las reorganiza según la causalidad del sistema: edge/cloud memory, expert parallel communication, throughput por lote, domain adaptation, RAG, modular experts, gradient flow y larger expert counts.

\subsection{Edge versus cloud: la ejecución dispersa no equivale a residencia dispersa}

En un memory-constrained edge device, un dense model suele ser la opción más directa, porque MoE sigue necesitando almacenar todos los experts. Si todos los pesos residen en la GPU, la resident memory se aproxima a los total parameters; si solo residen en la CPU, la GPU ahorra memoria, pero hay que transferirlos con frecuencia. El centralized serving cuenta con múltiples tarjetas, alto ancho de banda y lotes grandes, y por eso le resulta más fácil convertir la ejecución dispersa en ganancias de throughput.

\begin{importantbox}{La primera pregunta al elegir para edge}
Primero hay que preguntar si el dispositivo puede alojar los total weights, no solo cuántos active parameters hay. Un MoE con 12.9B activos aún puede requerir almacenar alrededor de 46.7B parámetros. Si una sola máquina solo puede alojar un dense de 7B, no se puede suponer que, por usar top-two, Mixtral ocupe también el equivalente a 13B.
\end{importantbox}

\teachervoice{00:31:45--00:33:10, Albert corrige a la persona que preguntó en la sala: en edge, un dense model puede ser más adecuado, porque la sparse inference sigue teniendo que cargar todos los experts; la ventaja de cost/performance de MoE corresponde mejor al serving centralizado.}

\subsection{Tamaño de lote y throughput: por qué un tráfico alto tiende a amortizar la complejidad}

Con lotes pequeños, los distintos experts reciben pocos tokens y de forma desigual, y es difícil amortizar el all-to-all y el kernel launch; un lote grande permite agrupar más tokens en expert GEMM de mayor tamaño, lo que eleva la arithmetic intensity. Puede representarse con un modelo simplificado de throughput:
\[
\operatorname{throughput}\approx \frac{T}{\max_i t_{\text{expert},i}+t_{\text{dispatch}}+t_{\text{combine}}}.
\]
Aumentar $T$ solo ayuda si los lotes de cada expert crecen y el imbalance se mantiene bajo control; si el expert más demandado aumenta su demanda en la misma proporción, la tail sigue determinada por él.

\begin{knowledgebox}{Leer métricas del sistema: la latency promedio no basta}
Reporte simultáneamente tokens/s, time-to-first-token, inter-token latency, P50/P95/P99, tamaño de lote por expert, all-to-all bytes y HBM occupancy. MoE puede aumentar el aggregate throughput y, al mismo tiempo, empeorar la latency de una solicitud individual o la tail.
\end{knowledgebox}

\teachervoice{00:42:47--00:43:33 y 00:54:56--00:56:17, el ponente enfatiza dos veces el use case: los lotes grandes y el high-volume serving muestran mejor el throughput de MoE; los entornos de pequeña escala y con memoria de GPU limitada cargan con más complejidad.}

\subsection{Communication topology: comunicarse entre nodes es más caro que entre GPU}

El costo del expert parallelism no depende solo de los bytes, sino también de la topology. Las GPU de un mismo node pueden comunicarse por NVLink/NVSwitch, mientras que entre nodes la comunicación pasa por InfiniBand/Ethernet y el network stack. Si un expert se distribuye a su vez en varias tarjetas con tensor parallel, el token dispatch y las collectives internas del expert se suman.

\begin{warningbox}{Más experts no siempre significa más velocidad}
Aumentar $E$ ofrece más routing choices, pero reduce el lote de cada expert, fragmenta más la placement y complica la metadata y la communication. Si los experts se reparten entre nodes, la latencia de red puede anular la ganancia del cómputo disperso.
\end{warningbox}

\teachervoice{00:40:43--00:41:36, Albert señala que la comunicación GPU-to-GPU ya es costosa y la node-to-node lo es aún más; el costo de comunicación depende del token routing y de la colocación en el hardware, y cómo seguir escalando sigue siendo un problema abierto.}

\subsection{Memory, offload y la dense-equivalent heuristic}

La sesión de Q\&A ofrece tres juicios que se confunden con facilidad. Primero, la naive resident memory se acerca a la proporción de pesos 46.7B/7B; segundo, el CPU offload ahorra HBM de la GPU, pero aumenta la transferencia de parámetros; tercero, estimar la dense-equivalent capability con $\sqrt{P_{\text{active}}P_{\text{total}}}$ solo puede ser una rule of thumb aproximada.

La heuristic de la media geométrica es
\[
P_{\text{equiv}}\approx \sqrt{P_{\text{active}}P_{\text{total}}},
\]
pero la capacidad también depende de los training tokens, la data quality, el optimizer, el router balance y la architecture. No sustituye un controlled scaling experiment con el mismo presupuesto de entrenamiento.

\teachervoice{00:50:42--00:52:42, el ponente dice que, cuando todo reside en la memoria de la GPU, la memory se acerca más a la total parameter ratio; el CPU offload pierde eficiencia de transferencia; la media geométrica sirve como rule of thumb, pero solo si los tokens de entrenamiento, la calidad y otras condiciones son comparables.}

\subsection{Knowledge versus reasoning: una benchmark category no es una sonda del mecanismo}

Alguien del público preguntó por qué Mixtral tiene mejor reasoning, y la respuesta de Albert fue mesurada: en los benchmarks de matemáticas, el modelo puede estar razonando de verdad o puede haber memorizado lemmas o patrones que luego invoca; la frontera entre knowledge y reasoning es difusa de por sí. Por eso, una puntuación más alta no demuestra que MoE haya aprendido un nuevo algoritmo interno, y mucho menos que attention/MLP se hayan separado por funciones.

\begin{warningbox}{Las etiquetas de capacidad no sustituyen la evidencia del proceso}
Para estudiar el reasoning mechanism hay que incorporar out-of-distribution composition, intermediate trace, counterfactual perturbation y memorization control. Basarse solo en que el nombre del benchmark contiene «reasoning» no permite determinar el proceso de cómputo.
\end{warningbox}

\teachervoice{00:53:26--00:54:41, el ponente considera que la ganancia de knowledge de Mixtral es clara, pero que la knowledge induzca mejoras en reasoning es muy speculative; los propios benchmarks difícilmente distinguen entre recall y reasoning.}

\subsection{Domain adaptation, RAG y modular experts son tres ejes independientes}

MoE aumenta la model capacity, el preentrenamiento continuado/fine-tuning cambia la domain distribution y RAG recupera evidencia externa durante la inference. Los tres pueden combinarse, pero no se sustituyen entre sí. Un medical dense model entrenado con datos especializados puede superar a un MoE general; un MoE también puede conectarse a RAG; insertar un nuevo medical expert en un checkpoint exige además entrenar el router para que aprenda cuándo usarlo.

\begin{knowledgebox}{Diferencias entre los tres ejes}
\begin{tabularx}{\textwidth}{L{0.18\textwidth} L{0.27\textwidth} X}
\toprule
Mecanismo & Qué cambia & Qué no resuelve automáticamente \\
\midrule
MoE & Capacidad de parámetros y token execution path & Actualización del conocimiento externo, domain label explícito \\
Domain tuning & Adaptación de los pesos a la distribución de datos objetivo & Hechos en tiempo real, provenance de la recuperación \\
RAG & Inference-time context/evidence & Capacidad de parámetros interna del modelo y routing balance \\
\bottomrule
\end{tabularx}
\end{knowledgebox}

\teachervoice{00:43:50--00:44:40, Albert dice que es difícil que un modelo general supere a uno con preentrenamiento/fine-tuning especializado de domain, y que los experts de Mixtral no son medical/coding experts explícitos. 00:59:02--00:59:43, explica además que RAG se combina de forma ortogonal con dense/MoE.}

\subsection{Expert swapping: que sea factible no significa que sea plug-and-play}

Si se sustituye un expert, los logits originales del router no corresponden al behavior del nuevo expert. Aunque la tensor shape coincida, el modelo no sabe cuándo enrutar los medical tokens al nuevo módulo, y además pueden alterarse las residual statistics. Como mínimo se necesita adaptation del router/nuevo expert, y verificar si la carga y la capacidad de los demás experts se degradan.

\begin{lstlisting}[caption={Flujo mínimo de validación tras sustituir un expert}]
load_base_moe_checkpoint()
replace_one_expert_with_domain_module()
freeze_or_partially_train_shared_attention()
train_router_and_new_expert_on_mixed_data()
measure_load_balance_and_general_capability()
run_domain_eval_and_out_of_domain_regression()
\end{lstlisting}

\teachervoice{00:59:43--01:00:54, el ponente considera que el expert swapping tiene potencial de investigación, pero exige explícitamente seguir entrenando las gating layers para que el router aprenda en qué momento debe invocarse el nuevo domain expert.}

\subsection{Gradient flow y training cost: \texorpdfstring{top-$k$}{top-k} sigue siendo optimizable de extremo a extremo}

Al final, la sesión de Q\&A corrige otro malentendido: aunque el router produce una selection top-$k$ discreta, las operaciones de gating y de los experts en la ruta seleccionada siguen siendo diferenciables, y el modelo puede entrenarse de extremo a extremo. En sentido estricto, la frontera de top-$k$ es por tramos/no suave, pero en la práctica la optimización propaga gradientes hacia los experts seleccionados en ese momento y se combina con load-balancing objectives.

El costo de entrenamiento puede resumirse como
\[
C_{\text{train}}\approx C_{\text{active forward/backward}}+C_{\text{router}}+C_{\text{communication}}+C_{\text{balance}}.
\]
Por lo general se acerca más al de un dense model del mismo active-size que al de un dense model del tamaño total (total-size), pero la extra communication y el almacenamiento de optimizer/state siguen presentes.

\begin{knowledgebox}{El optimizer state se sigue cobrando según los parámetros entrenados}
Adam/AdamW mantiene, para cada parámetro entrenable, el primer momento $m$ y el segundo momento $v$, además del parameter y el gradient. Si se entrenan todos los expert weights, el optimizer state depende de los total trainable parameters, no solo del active path del token actual; el entrenamiento distribuido suele requerir state sharding para reducir la memoria por tarjeta.
\end{knowledgebox}

\teachervoice{01:01:10--01:02:08, Albert dice que la routing path es differentiable de extremo a extremo y que el cómputo de entrenamiento se parece aproximadamente al de un active model de 13B, aunque además hay que asumir communication adicional.}

\subsection{Más experts: crece el espacio de elección y también la dificultad de servicio}

Aumentar el número de experts a 64/128 da al router un espacio de elección más fino y puede volver a los experts más especializados; pero el lote de cada expert se reduce, la placement abarca más nodes, el all-to-all se complica e incluso un solo expert podría no caber en una GPU. Los proveedores de modelos pueden ocultar esa complejidad con infraestructura especializada; los usuarios que despliegan por su cuenta tienen que enfrentarla.

\begin{importantbox}{Cuatro preguntas antes de aumentar $E$}
¿Cada expert es lo bastante grande para aprender features útiles? ¿Cuántos tokens recibe cada expert en cada step? ¿Los experts se reparten entre nodes? ¿El serving volume basta para formar GEMM grandes? Si estas preguntas no tienen respuesta afirmativa, más experts podrían solo aumentar el orchestration cost.
\end{importantbox}

\teachervoice{01:02:26--01:04:14, el ponente considera que más experts resultan muy atractivos para la investigación, pero vuelven más difíciles el multi-node serving y la communication; para una sola GPU, prefiere dense o un heavily quantized MoE, y juzga que Mixtral 8x7B cuantizado sigue siendo práctico.}

\subsection{Resumen de la sección}

La sesión de Q\&A convierte la exposición algorítmica de las slides en condiciones de despliegue: MoE es más adecuado para entornos con suficiente resident memory, interconnect y serving volume; en edge y con lotes pequeños no necesariamente compensa. Domain tuning, RAG y MoE son ortogonales entre sí; el expert swapping requiere router adaptation; el compute de entrenamiento se acerca al del active path, pero el optimizer state, la communication y el load balancing se siguen cobrando según el sistema más completo.
\section{Conclusión: la ganancia de la capacidad dispersa proviene de la selección, y su costo también}

La conclusión del ponente puede resumirse en tres juicios. Primero, SMoE intercambia activation dispersa por una mayor capacidad de conocimiento y puede situarse en una active-compute frontier excelente; segundo, el entrenamiento y el serving pueden ser eficientes, pero de ningún modo dependen solo del active parameter count; tercero, la expert specialization tiene estructura, pero no es tan directa como «experto en código/experto en matemáticas». La última slide deja abiertas como líneas de investigación tanto la architecture como la interpretability research.

\lecturefigure{slide-26-conclusion.jpg}{Síntesis de Albert Jiang sobre Sparse MoE, eficiencia y expert specialization}{Grabación oficial de Stanford Online 00:29:38}

\begin{importantbox}{Fórmula final de esta clase: las tres cuentas deben reportarse juntas}
\[
\boxed{\text{Model decision}=f(P_{\text{total}},P_{\text{active}},M_{\text{resident}},C_{\text{comm}},I_{\text{load}},Q_{\text{quality}})}
\]
Aquí, total/active describen la estructura; resident memory/communication/load describen el sistema, y quality describe la evidencia de la tarea. Cualquier comparación que conserve una sola variable puede inducir a error.
\end{importantbox}

\teachervoice{00:29:18--00:29:47, la síntesis original del ponente es: SMoE obtiene más knowledge gracias a la sparsity y puede entrenarse como un inference model eficiente, pero la expert specialization no es tan sencilla como sugiere la intuición, y en architecture e interpretability queda todavía mucho trabajo.}

\subsection{Resumen de la sección}

El valor de Mixtral no está en la etiqueta de «ocho expertos», sino en la token-level conditional computation: la capacidad total crece, la ruta activa se mantiene controlada y el router selecciona de forma dinámica. Al mismo tiempo, esto crea nuevos cuellos de botella de sistema y nuevas magnitudes observables para la interpretabilidad. Solo al considerar juntos la estructura, el hardware y la evidencia es posible juzgar si un MoE supera realmente a un dense baseline.
\section{Síntesis y ampliación}

Lo que más vale la pena trasladar de esta clase a otros modelos no es un conjunto de cifras de Mixtral, sino un método de auditoría: primero confirmar el dense baseline y luego describir con claridad el routing algorithm; reportar a la vez total/active/resident; incluir en el costo la communication y el load balance; y, por último, usar routing statistics e intervention para comprobar si los «experts» tienen de verdad una división de trabajo interpretable.

\subsection{Quince conclusiones centrales}

Las siguientes conclusiones se organizan en tres niveles: estructura, sistema e interpretación, para no tomar el nombre del modelo como respuesta.

\begin{enumerate}
  \item Mixtral conserva la attention de Mistral 7B y solo sustituye el MLP de cada capa por top-two-of-eight experts.
  \item GQA reduce los KV heads y la KV cache; SWA limita el attention neighborhood de cada capa.
  \item El Router puntúa cada token de forma independiente, elige el top-$k$, normaliza y combina con pesos los expert outputs.
  \item Cada capa tiene ocho experts; la numeración de los experts es permutable dentro de la capa y no tiene una identidad natural entre capas.
  \item Total parameters, active parameters y resident memory son tres magnitudes distintas.
  \item Los active parameters aproximan sobre todo el per-token compute; no representan la latency completa ni el dollar cost.
  \item El expert parallelism requiere un all-to-all de tokens, y la topology entre nodes puede convertirse en el cuello de botella.
  \item El load imbalance hace que el expert más ocupado determine la tail del step paralelo.
  \item Con lotes grandes y alto tráfico es más fácil agrandar los GEMM de cada expert y diluir el dispatch overhead.
  \item La MLP-knowledge hypothesis es sugerente desde el punto de vista experimental, pero no es un teorema de localización funcional.
  \item Los benchmarks de 2024 muestran que las mayores ganancias están en tareas knowledge-heavy, pero no prueban un nuevo reasoning algorithm.
  \item El discrete routing ofrece una señal que puede contabilizarse, pero la routing frequency no es una explicación semántica ni causal.
  \item Unos domain experts nítidos pueden dejar ociosos a otros experts y reducir el aprovechamiento.
  \item La expert ablation es una intervention importante, pero debe registrar la layer, la gate normalization y el tratamiento del overflow.
  \item MoE, domain tuning y RAG son ejes distintos y combinables; el expert swapping además requiere router adaptation.
\end{enumerate}

\subsection{Una tabla que enlaza estructura, entrenamiento y serving}

Para convertir los conceptos en revisiones de ingeniería, la tabla siguiente indica, para cada etapa, las variables que más conviene observar y los errores de juicio más comunes.

\begin{table}[H]
\centering
\small
\begin{tabularx}{\textwidth}{L{0.16\textwidth} L{0.27\textwidth} X}
\toprule
Etapa & Magnitudes clave observadas & Error de juicio común \\
\midrule
Architecture & $E,k$, shared/expert params, GQA/SWA & Estimar el número de parámetros a partir del nombre del modelo \\
Training & tokens/expert, aux loss, gradient, optimizer state & Calcular la memoria solo con los active params \\
Placement & expert sharding, node topology, HBM & Suponer que los experts no activados no necesitan residir en memoria \\
Serving & lote, all-to-all, imbalance, tail latency & Usar FLOPs en lugar del cost real \\
Interpretation & routing stats, ablation, counterfactual & Ponerle directamente una etiqueta de domain a cada expert \\
Evaluation & task setup, ejes active/total, history date & Tomar el plot de 2024 como ranking actual \\
\bottomrule
\end{tabularx}
\caption{Marco mínimo de auditoría de MoE, desde la arquitectura hasta la interpretación.}
\end{table}

\subsection{MoE release checklist para la práctica}

La siguiente checklist sirve para revisar artículos, model cards o releases internos; su objetivo principal es evitar que solo se muestre el parameter headline.

\begin{lstlisting}[caption={Sparse MoE release checklist}]
report_total_and_active_parameters()
report_expert_count_topk_and_shared_modules()
report_weight_precision_and_resident_memory()
report_tokens_per_expert_and_load_imbalance()
report_expert_parallel_topology_and_all_to_all_cost()
compare_against_dense_models_at_compute_and_memory_budgets()
publish_routing_analysis_with_causal_controls()
document_quantization_offload_and_failure_cases()
\end{lstlisting}

\begin{warningbox}{Al ver «13B active», pregunte de inmediato}
¿Cuántos parámetros totales tiene el checkpoint? ¿Dónde están alojados todos los experts: en una sola tarjeta, en un solo node o repartidos entre nodes? ¿De qué tamaño es el lote? ¿Cuál es la P99 latency? ¿Hay token drop/reroute? Sin esta información, el active parameter headline solo describe una ruta de cómputo parcial.
\end{warningbox}

\subsection{Preguntas de autoevaluación}

Estas preguntas sirven para comprobar si de verdad se dominan las tres cuentas y la routing evidence.

\begin{enumerate}
  \item ¿Por qué Mixtral no es un ensemble de ocho modelos 7B independientes?
  \item ¿Qué recurso reduce GQA y cuál reduce SWA?
  \item Escriba el proceso del top-two routing, desde los logits hasta la weighted sum.
  \item ¿Por qué $8\times 7\mathrm{B}$ no es el valor exacto de los total ni de los active parameters?
  \item ¿Por qué los active parameters no permiten predecir directamente la serving latency?
  \item ¿Qué hace el expert parallel all-to-all en cada una de las dos fases, dispatch y combine?
  \item ¿Por qué el expert más lento controla el tiempo de finalización de un lote?
  \item ¿En qué se diferencian quantization, offload y sparsification?
  \item ¿Qué puede probar un domain routing histogram y qué no?
  \item ¿Por qué unos experts dedicados solo a código pueden reducir el aprovechamiento global?
  \item ¿Qué variables de control debe registrar una expert ablation?
  \item ¿Por qué RAG y MoE no son alternativas entre sí?
  \item ¿Por qué, tras sustituir un expert, todavía hay que entrenar el router?
  \item ¿En qué condiciones más experts podrían, por el contrario, reducir el throughput?
\end{enumerate}

\subsection{Lecturas adicionales}

\begin{itemize}
  \item Jiang et al., \href{https://arxiv.org/abs/2401.04088}{\emph{Mixtral of Experts}}: artículo principal sobre la estructura del modelo, su desempeño y el routing analysis.
  \item Jiang et al., \href{https://arxiv.org/abs/2310.06825}{\emph{Mistral 7B}}: GQA, SWA y el dense baseline.
  \item Shazeer et al., \href{https://arxiv.org/abs/1701.06538}{\emph{Outrageously Large Neural Networks}}: la sparsely-gated MoE layer.
  \item Fedus et al., \href{https://arxiv.org/abs/2101.03961}{\emph{Switch Transformers}}: entrenamiento disperso a gran escala y parallelism.
  \item Al leer estos artículos, registre también total/active parameters, precision, hardware, lote y communication; no se limite a extraer el benchmark headline.
\end{itemize}

\end{document}
